空间转动改变坐标分量却保持欧氏距离，量子态的整体相位变换改变复相位却保持全部概率；这两件熟悉的事实都具有同一结构：允许的变换可以连续复合，而且复合后物理不变量保持不变。群论把这种“变换怎样复合”抽象出来。Lorentz 群、$U(1)$、$SU(2)$ 和 $SU(3)$ 的差别在于各自保持什么对象不变，但它们都服从同一组群运算规则。

先设 $G$ 是一组可以依次执行的变换，并把“先做 $g_2$、再做 $g_1$”记成乘积 $g_1g_2$。若这组变换满足
\begin{equation}
 g_1g_2\in G,\qquad
 (g_1g_2)g_3=g_1(g_2g_3),\qquad
 eg=ge=g,\qquad
 gg^{-1}=g^{-1}g=e,
 \label{eq:group-axioms-math-dp76}
\end{equation}
就称 $G$ 为一个\emph{群}。四个条件依次表示：两个允许变换连续做仍是允许变换；连续做三次时括号位置不影响结果；存在“什么也不做”的单位变换 $e$；每个变换都存在可以把它撤销的逆变换。群不要求乘法可交换，所以一般可以有 $g_1g_2\neq g_2g_1$。例如三维空间绕不同轴的有限转动通常不交换，这已经是非交换群最直观的例子。

当群元素可以由若干连续实参数平滑标记时，离散的复合规则还能进一步被微分。严格地说，若群元素构成光滑流形，而且群乘法与取逆都是光滑映射，就称为\emph{Lie 群}。这里“流形”只需理解成：每个群元素附近都能选取若干实数作为局部坐标，并对穿过该点的群曲线像普通多元函数那样求导。$SO(1,3)$、$U(1)$ 和 $SU(N)$ 都可表示为满足光滑矩阵约束的有限维矩阵集合，因此其无穷小变换可以直接用矩阵微分处理。

Lie 群最重要的简化发生在单位元附近。取一条通过单位元的光滑群曲线 $g(t)$，满足 $g(0)=e$。它在 $t=0$ 的速度
\begin{equation*}
 X\equiv\left.\frac{\dd g(t)}{\dd t}\right|_{t=0}
\end{equation*}
描述“沿某个方向做无穷小变换”。所有这类速度组成一个实向量空间，称为单位元处的\emph{切空间}：
\begin{equation}
 \boxed{\mathfrak g\equiv T_eG.}
 \label{eq:lie-algebra-tangent-general-dp68}
\end{equation}
这个向量空间连同其 Lie 括号称为群 $G$ 的\emph{Lie 代数}。它不是另一套独立对称性，而是有限群变换在单位元附近的一阶微分结构。选择一组基 $\{X_a\}_{a=1}^{d_G}$，就等于选择 $d_G$ 个独立的无穷小变换方向，其中
\begin{equation}
 d_G\equiv\dim_{\mathbb R}\mathfrak g
 \label{eq:lie-algebra-dimension-dp68}
\end{equation}
是这些独立方向的数目。于是单位元附近的一般变换可写成
\begin{equation}
 g(\boldsymbol\theta)
 =I+\theta^aX_a+\order(\theta^2).
 \label{eq:lie-infinitesimal-dp11}
\end{equation}
这里的 $X_a$ 是群流形单位元处的切向量；当群进一步表示为量子 Hilbert 空间上的幺正算符时，相应切向表示才给出熟悉的厄米量子生成元。

若只沿一个固定方向 $X$ 连续变换，就得到单参数子群 $g(t)$，它满足
\begin{equation}
 g(t_1)g(t_2)=g(t_1+t_2),
 \qquad g(0)=I.
 \label{eq:one-param-group-dp11}
\end{equation}
对矩阵 Lie 群，这个函数由初始速度 $X=\dot g(0)$ 唯一确定为
\begin{equation}
 \boxed{g(t)=\exp(tX),}
 \label{eq:one-param-exp-dp11}
\end{equation}
其中矩阵指数就是普通幂级数
\begin{equation}
 e^X\equiv\sum_{n=0}^{\infty}\frac{X^n}{n!}.
 \label{eq:matrix-exponential-series-dp68}
\end{equation}
因此“生成元”这个词的含义很具体：知道单位元处的无穷小方向 $X$，指数映射就沿这个方向生成有限连续变换。

连续群还必须记录不同无穷小方向的先后次序是否重要。对所用的矩阵 Lie 群，群的共轭映射会把单位元切空间映回自身，因此两个生成方向的矩阵对易子仍是合法的无穷小生成方向。矩阵 Lie 群的 Lie 括号于是为
\begin{equation}
 \boxed{[X,Y]_{\mathfrak g}=XY-YX,\qquad X,Y\in\mathfrak g.}
 \label{eq:matrix-lie-bracket-math-dp76}
\end{equation}
它测量两个无穷小变换方向的非交换部分；若该括号恒为零，相应 Lie 代数就是阿贝尔的。

在基 $X_a$ 上展开后，必存在一组实数 $f_{ab}{}^c$ 使
\begin{equation}
 \boxed{[X_a,X_b]_{\mathfrak g}=f_{ab}{}^cX_c.}
 \label{eq:lie-algebra-general-dp11}
\end{equation}
这些数称为\emph{结构常数}：它们记录“先沿第 $a$ 个无穷小方向、再沿第 $b$ 个方向”所产生的非交换效应如何重新分解到原来的生成方向上。矩阵乘法的结合律使对易子满足 Jacobi 恒等式
\begin{equation}
 [A,[B,C]]+[B,[C,A]]+[C,[A,B]]=0,
 \label{eq:commutator-jacobi-dp68}
\end{equation}
因而
\begin{equation}
 \boxed{
 f_{bc}{}^{d}f_{ad}{}^{e}
 +f_{ca}{}^{d}f_{bd}{}^{e}
 +f_{ab}{}^{d}f_{cd}{}^{e}=0.}
 \label{eq:jacobi-structure-dp11}
\end{equation}
Baker--Campbell--Hausdorff 公式只是把同一件事提升到有限但仍靠近单位元的变换：
\begin{equation}
 e^Xe^Y
 =\exp\!\left[
 X+Y+\frac12[X,Y]
 +\frac1{12}[X,[X,Y]]+\frac1{12}[Y,[Y,X]]+\cdots
 \right].
 \label{eq:bch-lie-dp11}
\end{equation}
它说明有限变换的非交换修正由同一个 Lie 括号逐级组织。

抽象群只规定变换怎样彼此复合；物理计算还需要知道同一个变换怎样作用在具体对象上。例如空间转动既可以作用在三维位置向量上，也可以作用在某个具有多个复分量的量子态向量上；两者使用的矩阵维数可以不同，却描述同一个抽象转动。\emph{表示}就是把这种“同一抽象变换在某个向量空间上的具体作用”数学化。令 $V_R$ 是一组物理分量所在的复向量空间，$GL(V_R)$ 表示 $V_R$ 上所有可逆线性变换组成的群。一个群表示是映射
\begin{equation*}
 D_R:G\longrightarrow GL(V_R)
\end{equation*}
并满足
\begin{equation*}
 D_R(g_1g_2)=D_R(g_1)D_R(g_2),
 \qquad D_R(e)=I_R.
\end{equation*}
第一条条件保证“先后做两个抽象对称变换”与“把对应矩阵先后作用在物理分量上”完全一致。数学上，把乘法关系保持不变的映射称为\emph{群同态}；因此群表示就是一个取值在可逆线性变换群中的群同态。表示空间的复维数记为
\begin{equation*}
 d_R\equiv\dim_{\mathbb C}V_R.
\end{equation*}
如果存在一个既非零也非整个 $V_R$ 的子空间，并且所有 $D_R(g)$ 都把它映回自身，那么表示还能继续分块；若不存在这样的真不变子空间，就称该表示为\emph{不可约表示}。不可约表示是对称性能够区分的最小独立块，所以后面“粒子按质量、自旋分类”最终会落到寻找 Poincar\'e 群的不可约表示。

量子内部对称通常要求保持态的内积，因此采用酉表示。此时表示在单位元附近的导数 $dD_R(X_a)$ 是反厄米算符。为了使用量子力学熟悉的厄米可观测量 convention，定义
\begin{equation}
 \boxed{T_R^a\equiv \ii\,dD_R(X_a),}
 \qquad
 D_R(\ee^{\theta^aX_a})
 =\exp(-\ii\theta^aT_R^a).
 \label{eq:unitary-representation-generators-dp68}
\end{equation}
因为表示保持群乘法，它的微分也保持 Lie 括号：
\begin{equation*}
 dD_R([X_a,X_b])
 =[dD_R(X_a),dD_R(X_b)].
\end{equation*}
于是具体表示中的厄米生成元满足
\begin{equation}
 \boxed{[T_R^a,T_R^b]=\ii f_{ab}{}^{c}T_R^c.}
 \label{eq:rep-lie-algebra-dp11}
\end{equation}
这里要区分三个层级：$X_a$ 是抽象群的无穷小方向，$f_{ab}{}^c$ 描述群本身的局部乘法结构，而 $T_R^a$ 是这些方向在特定表示空间 $V_R$ 上的矩阵实现。改变表示会改变矩阵 $T_R^a$，但不会改变抽象群是哪一个群。



生成元的对易关系描述群的局部乘法结构，但不能单独区分不同表示的尺度与二次不变量。为此需要基底无关的表示不变量。群若作为参数空间是紧致的，直观上表示其连续参数不会沿某个方向无限逃逸；$U(N)$ 与 $SU(N)$ 是典型紧致矩阵群。Lie 代数若除了 $\{0\}$ 和整个代数外，不存在对所有外部生成元作 Lie 括号后仍封闭的非零真子空间，就称为简单 Lie 代数；这样的封闭子空间称为理想。$SU(2)$ 与 $SU(3)$ 的 Lie 代数都是紧致简单 Lie 代数，因此其不可约表示可以用一组标准二次不变量比较。

在紧致简单 Lie 代数上，可以选择一个被全部群变换保持的不变内积，并把基归一化到该内积的分量为 $\delta^{ab}$。这相当于为“生成元方向之间怎样取长度和夹角”固定统一标尺。在不可约酉表示 $R$ 中，生成元的二次迹只能与这个不变内积成正比：
\begin{equation}
 \boxed{
 \Tr_R(T_R^aT_R^b)=T(R)\delta^{ab}.}
 \label{eq:dynkin-index-dp11}
\end{equation}
比例系数 $T(R)$ 衡量表示中生成元的整体迹归一化；它通常称为表示 $R$ 的 Dynkin 指标。它不是脱离生成元归一化的绝对数值：若所有生成元整体重标度，$T(R)$ 与二次 Casimir $C_2(R)$ 会同步改变。

同一组生成元还定义算符
\begin{equation}
 \mathcal C_R\equiv\sum_{a=1}^{d_G}T_R^aT_R^a.
 \label{eq:quadratic-invariant-operator-dp68}
\end{equation}
利用结构常数对生成元指标的反对称性可得 $[\mathcal C_R,T_R^b]=0$。这里使用一个基本线性代数事实：若一个复不可约表示中的线性算符与所有表示矩阵都对易，那么它不能在表示空间内再挑出新的不变方向，只能是恒等算符的常数倍。这就是 Schur 引理。因此
\begin{equation}
 \boxed{
 \mathcal C_R=C_2(R)I_{d_R}.}
 \label{eq:casimir-general-rep-dp11}
\end{equation}
这个与全部生成元对易的二次算符通常称为二次 Casimir 算符，$C_2(R)$ 是它在表示 $R$ 上的本征值。对可约表示，应在每个不可约块上分别使用这一式。

现在对\eqnrefs{eq:casimir-general-rep-dp11} 取整个表示空间的迹，并用\eqnrefs{eq:dynkin-index-dp11}：
\begin{align}
 d_RC_2(R)
 &=\sum_{a=1}^{d_G}\Tr_R(T_R^aT_R^a),\\
 &=T(R)\sum_{a=1}^{d_G}\delta^{aa}
 =d_GT(R).
\end{align}
于是得到一般恒等式
\begin{equation}
 \boxed{d_RC_2(R)=d_GT(R).}
 \label{eq:casimir-dynkin-relation-dp11}
\end{equation}
这里三个对象必须区分：$d_G=\dim_{\mathbb R}\mathfrak g$ 是该简单 Lie 代数的生成元数目，$d_R=\dim_{\mathbb C}V_R$ 是表示空间维数，而 $T(R)$ 与 $C_2(R)$ 是在既定生成元归一化下定义的表示不变量。\eqnrefs{eq:casimir-dynkin-relation-dp11} 是后面所有 $C_F,C_A,T_F$ 数值的母关系，而不是由这些特殊数值反推的一般化。

最简单的紧 Lie 群 $U(1)$ 只有一个生成方向，因此 Lie 代数是阿贝尔的。由于阿贝尔生成元可以与耦合常数作相反的整体重标度，$U(1)$ 本身没有由非阿贝尔结构固定的唯一归一化。在一维不可约酉表示中，生成元只是实数 $y$，可写成
\begin{equation}
 \psi\longmapsto \ee^{-\ii y\theta}\psi.
 \label{eq:u1-charge-rep-dp11}
\end{equation}
所以阿贝尔“荷”在数学上就是选定生成元归一化后的一维表示标签。具体理论可以把这个实数解释成不同的守恒荷；它究竟对应哪种物理电荷，必须等作用量和相互作用建立后再由理论本身确定。

现在把抽象结构代入最常见的矩阵群。$U(N)$ 是所有 $N\times N$ 酉矩阵组成的群，即 $U^\dagger U=I_N$；其中再要求行列式为 $1$ 的子群称为 $SU(N)$，字母 $S$ 表示 ``special''。最小的非阿贝尔例子是 $SU(2)$：

$SU(2)$ 定义为
\begin{equation*}
 SU(2)=\{U\in\mathbb C^{2\times2}\mid U^\dagger U=I_2,\ \det U=1\}.
\end{equation*}
所谓\emph{基本表示}，就是让群矩阵直接作用在定义它的列向量空间 $\mathbb C^2$ 上；它的维数为 $2$。在这一本征的二维表示中，可取厄米生成元 $T^a=\sigma^a/2$。Pauli 矩阵满足
\begin{equation}
 \sigma^a\sigma^b
 =\delta^{ab}I_2+\ii\epsilon^{abc}\sigma^c,
 \label{eq:pauli-product-lie-dp68}
\end{equation}
因此
\begin{equation}
 \boxed{[T^a,T^b]=\ii\epsilon^{abc}T^c.}
 \label{eq:su2-algebra-dp11}
\end{equation}
定义 $T^\pm=T^1\pm\ii T^2$ 后，
\begin{equation}
 [T^3,T^\pm]=\pm T^\pm,
 \qquad [T^+,T^-]=2T^3.
 \label{eq:su2-ladder-algebra-dp11}
\end{equation}
一般不可约表示可由半整数 $j$ 标记，其二次 Casimir 与 $T^3$ 本征值为
\begin{equation}
 \boxed{
 C_2(j)=j(j+1),
 \qquad m=-j,-j+1,\ldots,j.}
 \label{eq:su2-casimir-dp11}
\end{equation}
$j=1/2$ 给二维基本表示，$j=1$ 给三维表示。这里的 $j$ 只标记抽象 $SU(2)$ 表示；两个具有同构 Lie 代数的 $SU(2)$ 对称完全可以作用在不同的向量空间上，因此表示标签本身还不指定具体物理含义。

同理，对一般 $SU(N)$，基本表示就是矩阵 $U$ 直接作用在 $\mathbb C^N$ 上，所以表示空间维数为 $d_F=N$。群还可以作用在自己的无穷小生成方向上：给定一个 Lie 代数元素 $X$，有限群元素 $U$ 通过共轭变换 $X\mapsto UXU^{-1}$ 把它线性混合成其它生成方向。因为被变换的向量空间就是 Lie 代数本身，这个表示称为\emph{伴随表示}。对 $SU(2)$，上面的 $j=1$ 三维表示正是这一伴随表示。一般 $SU(N)$ 的 Lie 代数有 $d_G$ 个独立生成方向，因此伴随表示的维数也是
\begin{equation*}
 d_G=N^2-1.
\end{equation*}
粒子物理通常把基本表示生成元的归一化约定固定为
\begin{equation}
 \Tr_F(T^aT^b)=T_F\delta^{ab},
 \qquad \boxed{T_F\equiv T(F)=\frac12.}
 \label{eq:generator-trace-normalization-dp11}
\end{equation}
这里 $T_F=1/2$ 是生成元归一化的选择，而不是先于一般表示理论的独立动力学结论。代入一般恒等\eqnrefs{eq:casimir-dynkin-relation-dp11} 立即得到基本表示的二次 Casimir
\begin{equation}
 \boxed{C_F\equiv C_2(F)=\frac{N^2-1}{2N}.}
 \label{eq:cf-general-math-dp11}
\end{equation}
同一归一化下，伴随表示生成元可写成 $(F^a)_{bc}=-\ii f^{abc}$；$SU(N)$ 的伴随二次 Casimir 为
\begin{equation}
 \boxed{C_A\equiv C_2(\mathrm{Adj})=N.}
 \label{eq:ca-general-math-dp11}
\end{equation}
基本生成元还满足完备关系
\begin{equation}
 \boxed{
 \sum_{a=1}^{N^2-1}(T^a)_{ij}(T^a)_{kl}
 =T_F\left(\delta_{il}\delta_{jk}-\frac1N\delta_{ij}\delta_{kl}\right).}
 \label{eq:suN-fierz-math-dp68}
\end{equation}
它只属于 $SU(N)$ 基本表示这一特例，用于把基本表示生成元的求和化为 Kronecker 张量。作为一个常用的 $N=3$ 特例，有
\begin{equation*}
 T_F=\frac12,\qquad C_F=\frac43,\qquad C_A=3.
\end{equation*}
因此这些数值只是一般 $SU(N)$ 表示结构在 $N=3$ 的基本表示与伴随表示上的末端代入。

相互作用项同时含多个场，因此还需要张量积表示。若 $v\in V_R$、$w\in V_S$，则
\begin{equation}
 D_{R\otimes S}(g)=D_R(g)\otimes D_S(g),
 \label{eq:tensor-product-rep-dp11}
\end{equation}
其无穷小生成元为
\begin{equation}
 \boxed{
 T_{R\otimes S}^a
 =T_R^a\otimes I+I\otimes T_S^a.}
 \label{eq:tensor-product-generator-dp11}
\end{equation}
若张量积包含一维平凡表示，就存在不变张量把各表示指标收缩成群不变量。局域多场乘积能够进入对称作用量的条件，正是其表示指标能够借助这类不变张量收缩成单态。

对任意 $2\times2$ 矩阵 $U$，完全反对称张量
\begin{equation*}
 \epsilon=\begin{pmatrix}0&1\\-1&0\end{pmatrix}=\ii\sigma^2
\end{equation*}
满足
\begin{equation}
 U^T\epsilon U=(\det U)\epsilon.
 \label{eq:2x2-epsilon-det-dp68}
\end{equation}
因此对 $U\in SU(2)$，
\begin{equation}
 \boxed{U^T\epsilon U=\epsilon.}
 \label{eq:su2-epsilon-invariance-dp11}
\end{equation}
这意味着 $u^T\epsilon v$ 是 $SU(2)$ 单态。共轭二重态定义为
\begin{equation}
 \boxed{\widetilde\Phi\equiv\ii\sigma^2\Phi^*,}
 \label{eq:tilde-higgs-math-dp11}
\end{equation}
并仍按基本二重态变换。这个构造只依赖 $SU(2)$ 的不变反对称张量，后面遇到需要把共轭场重新组织成二重态的物理模型时可以直接使用。

对 $SU(N)$ 的基本表示与反基本表示，$\delta_i{}^j$ 是不变张量。三维完全反对称张量与行列式满足
\begin{equation}
 U_i{}^{i'}U_j{}^{j'}U_k{}^{k'}\epsilon_{i'j'k'}
 =(\det U)\epsilon_{ijk}.
 \label{eq:determinant-epsilon-identity-dp68}
\end{equation}
因此对 $SU(3)$，
\begin{equation}
 \boxed{
 U_i{}^{i'}U_j{}^{j'}U_k{}^{k'}\epsilon_{i'j'k'}
 =\epsilon_{ijk}.}
 \label{eq:su3-epsilon-invariant-dp11}
\end{equation}
三个基本三维表示因而可以通过 $\epsilon_{ijk}$ 收缩成群单态。

除二次迹外，还可以从三个生成元构造完全对称的表示不变量
\begin{equation}
 d_R^{abc}
 \equiv2\Tr_R\!\left[T_R^a\{T_R^b,T_R^c\}\right].
 \label{eq:anomaly-symmetric-trace-dp11}
\end{equation}
它只由表示 $R$ 和生成元归一化决定。先把它作为群表示的三阶对称张量保存下来；它在量子理论中的具体用途等相应圈图与量子对称性问题建立后再解释。

Lie 群本身的局部结构由单位元切空间与 Lie 括号决定；表示把同一抽象 Lie 代数映到具体内部向量空间；$T(R)$ 与 $C_2(R)$ 是在既定生成元归一化下定义的二次表示不变量，并先满足 $d_RC_2(R)=d_GT(R)$，随后才有 $SU(N)$ 基本表示中的 $T_F,C_F,C_A$。张量积与不变张量决定多场局域项能否组成群单态。
